
%%%%%%%%%%%%%%%%%%%%%%%%%%%%%%%%%%%%%%%%%%%%%%%%%
%
%
%                      Noms et Numéros de cours
%
%
%%%%%%%%%%%%%%%%%%%%%%%%%%%%%%%%%%%%%%%%%%%%%%%%%


\COURS{07}{machinesdeTuring}{Turing machines}
\COURS{07}{modelsofcomputation}{A few other models of computation}


%====================
%                      INF412, CSE304, LFMI2020
%=====================

\def\POLYINF412{
	\def\XPOLY{}
	\providecommand\TITRECOURS{\LANGUE{Foundations of Computer Science}{Fondements de l'informatique} \\%[3cm] % \\
                                % INF412
\LANGUE{Logic, models, and computations}{Logique, mod\`eles, et calculs}}
	\providecommand\SUBTITRECOURS{\LANGUE{Course}{Cours}  CSC\_41012\_EP \\[0.2cm] 
	\LANGUE{of}{de} l'Ecole  Polytechnique %\\ \LANGUE{(previous course code: INF412}{(ancien code du cours: INF412)}
}}

\def\POLYCSE304{
	\def\XPOLY{}
	\providecommand\TITRECOURS{Computability and Complexity.}
	\providecommand\SUBTITRECOURS{Course CSE304. \\%\\[0.5cm]  
	Bachelor Ecole Polytechnique.}
}

\def\POLYLFMIvingtvingt{
	%\def\XPOLY{}
	%\providecommand\TITRECOURS{Computability and Complexity.}
	%\providecommand\SUBTITRECOURS{Course LMFI. \\%\\[0.5cm]  
	%Bachelor Ecole Polytechnique.}
}

\def\POLYMPRIvingtvingtun{
	%\def\XPOLY{}
	%\providecommand\TITRECOURS{Computability and Complexity.}
	%\providecommand\SUBTITRECOURS{Course LMFI. \\%\\[0.5cm]  
	%Bachelor Ecole Polytechnique.}
}

%====================
%                      INF412
%=====================

\INFquatrecentdouze{00}{polyinfquatrecentdouze}{Logique, Modèles de
Calculs, Complexité}
\INFquatrecentdouze{01}{introductionquatrecentdouze}{Introduction}
\INFquatrecentdouze{02}{preliminairesquatrecentdouze}{Récursivité et induction}
\INFquatrecentdouze{03}{propositionnel}{Calcul propositionnel}
\INFquatrecentdouze{04}{preuvespropositionnelles}{Démonstrations}
\INFquatrecentdouze{05}{premierordre}{Calcul des prédicats}
\INFquatrecentdouze{06}{completude}{Modèles. Complétude.}
\INFquatrecentdouze{07}{machinesdeTuringquatrecentdouze}{Machines de Turing}
\INFquatrecentdouze{07}{modelesdecalculsquatrecentdouze}{Autres modèles de calculs}
\INFquatrecentdouze{08}{algorithmes}{Qu'est-ce qu'un algorithme
séquentiel? Thèse de Church.}
\INFquatrecentdouze{09}{calculabilitequatrecentdouze}{Calculabilité}
\INFquatrecentdouze{10}{incompletude}{Incomplétude de l'arithmétique}
\INFquatrecentdouze{11}{basescomplexite}{Bases de l'analyse de complexité d'algorithmes}
\INFquatrecentdouze{12}{complexitetempsquatrecentdouze}{Complexité en temps}
\INFquatrecentdouze{13}{circuitsquatrecentdouze}{La notion de circuit}
\INFquatrecentdouze{14}{npcompletude}{Quelques problèmes
$\NP$-complets}
\INFquatrecentdouze{14}{npcompletudeenglish}{Some $\NP$-complete
problems}
\INFquatrecentdouze{15}{autresclassescomplexite}{Space complexity}
\INFquatrecentdouze{15}{autresclassescomplexitequatrecentdouze}{Complexité en espace mémoire}
\INFquatrecentdouze{16}{corrections}{Solutions de certains exercices}
\INFquatrecentdouze{17}{tvracannexeinfquatrecentdouze}{Vrac annexe}


%====================
%                      Logique
%=====================
  
\COURS{01}{theoriedesensemblesZFC}{Théorie des ensembles ZFC,
  Hiérarchie de Von Neumann, Hiérarchie de Gödel / des constructibles}

\COURS{02}{ordinaux}{Ordinals}

\COURS{03}{admissibilite}{Ensembles admissibles, Théorie des
  ensembles de KP}

\COURS{30}{theoriedesmodeles}{Théorie des modèles. Logique}

\COURS{06}{definissabilite}{Definissability. Coding}


\COURS{05}{reels}{Reals}

\COURS{10}{incompletudeenglish}{Incompleteness of arithmetic}

%====================
%                      Algo
%=====================

%% VERSION Algo maitrise Nancy

\NANCY{01}{approximation-05.ps@}{approximation-05.ps@}
\NANCY{02}{progdynamique-trianglepascal.ps@}{progdynamique-trianglepascal.ps@}
\NANCY{03}{motivation-complexite-03.ps@}{motivation-complexite-03.ps@}
\NANCY{04}{progdynamique-trianglepascal2.ps@}{progdynamique-trianglepascal2.ps@
}
\NANCY{05}{algorithmes-glouton-02.ps@}{algorithmes-glouton-02.ps@}
\NANCY{06}{np-completude-04.ps@}{np-completude-04.ps@}

	

%====================
%                      Verification
%=====================

%% VERSION M2 2005

\NANCY{01}{01-N-Generalites-Rappels}{01-N-Generalites-Rappels}

\NANCY{01p5}{01p5-N-Explosion-Combinatoire}{01p5-N-Explosion-Combinatoire}

\NANCY{02}{02-N-Minimisation-Graphe}{02-N-Minimisation-Graphe}

\NANCY{03}{03-N-Symetries}{03-N-Symetries}

\NANCY{04}{04-N-Cone-Influence.ps@}{04-N-Cone-Influence.ps@}

\NANCY{05}{
  05-N-Raffinement-Automatique-Partition.ps@}{05-N-Raffinement-Automatique-Partition.ps@}

\NANCY{06}{06-N-Pouvoir-Distinction-Logiques.ps@}{
  06-N-Pouvoir-Distinction-Logiques.ps@}

\NANCY{07}{  07-N-PreordreStructures.ps@}{
  07-N-PreordreStructures.ps@}

\NANCY{08}{
  08-N-Raisonnement-Compositionel.ps@}{08-N-Raisonnement-Compositionel.ps@}

\NANCY{09}{ 09-N-Reduction-OrdrePartiel.ps@}{
  09-N-Reduction-OrdrePartiel.ps@}

\NANCY{10}{10-N-AutomatesTemporises.ps@}{10-N-AutomatesTemporises.ps@}

\NANCY{11}{11-N-MatricesDifferencesBornees.ps@}{11-N-MatricesDifferencesBornees.ps@}

\NANCY{12}{12-N-SystemesHybrides.ps@}{12-N-SystemesHybrides.ps@}

\NANCY{13}{13-N-ThBuchi.ps@}{13-N-ThBuchi.ps@}

\NANCY{14}{14-N-Verif-Probabiliste.ps@}

\NANCY{xA}{xA-TheorieLangagesTemporises.ps@}{xA-TheorieLangagesTemporises.ps@}

\NANCY{xA}{xA-N-Explosion-Combinatoire.ps@}{xA-N-Explosion-Combinatoire.ps@}

\NANCY{xB}{xB-TheorieOmegaAutomates.ps@}{xB-TheorieOmegaAutomates.ps@}

\NANCY{xB}{xB-N-Theorie-Des-Logiques-Et-Automates.ps@}{xB-N-Theorie-Des-Logiques-Et-Automates.ps@}

\NANCY{xC}{xC-AutomatesRectangulaires.ps}{xC-AutomatesRectangulaires.ps}

\NANCY{xD}{xD-TerminaisonAlgos.ps@}{xD-TerminaisonAlgos.ps@}

\NANCY{xE}{ xE-ModelChecking-CTLe.ps@}{ xE-ModelChecking-CTLe.ps@}

\NANCY{xF}{xF-TheorieLangagesTemporises.ps@}

\NANCY{xZ}{ xZ-Reste.ps@}{ xZ-Reste.ps@}
 



%====================
%                      Descriptive set theory
%=====================


%%
\COURS{04}{trees}{
  Trees}

\COURS{10-A}{dstbaire}{ Descriptive Set Theory: Baire space}

\COURS{10-B}{dstborelhierarchy}{ Descriptive Set Theory: Borel
  Hierarchy}

\COURS{10-C}{dsteffectivehierarchy}{ Descriptive Set Theory: Effective
  Hierarchy}

\COURS{10-D}{dstanalytichierarchy}{ Descriptive Set Theory: Analytic
  Hierarchy}

\COURS{11}{hierarchykechriswoodin}{ Kechris and Woodin's hierarchy}


\COURS{12}{complexitedumonde}{Complexity
  In Analysis}
%%

%========================================
%                      Complexité / Classique
%========================================

\NANCY{00}{00-N-Motivation-Theorie-Complexite.pdf@}{00-N-Motivation-Theorie-Complexite.pdf@}

\NANCY{01}{01-N-Machines-De-Turing.ps@}{01-N-Machines-De-Turing.ps@}

\NANCY{02}{02-N-These-De-Church-And-Co.pdf@}{02-N-These-De-Church-And-Co.pdf@}

\NANCY{03}{03-N-Classes-Complexite-Standards.pdf@}{03-N-Classes-Complexite-Standards.pdf@}

\NANCY{04}{04-N-Problemes-Complets.pdf@}{04-N-Problemes-Complets.pdf@}

\NANCY{05}{05-Classes-Paralleles.pdf@ }{05-Classes-Paralleles.pdf@ }

\NANCY{06}{06-Classes-Probabilistes.pdf@}{06-Classes-Probabilistes.pdf}

\NANCY{07}{07-Caracterisations-Logiques.pdf@}{07-Caracterisations-Logiques.pdf@}

\NANCY{08}{08-N-Comptage-Inductif.pdf@}{08-N-Comptage-Inductif.pdf@}


\INFcinqsixun{30}{polycomplexite}{Polycopié Complexité}
\INFcinqsixun{30}{introductioncomplexite}{Introduction à  la complexité}
\INFcinqsixun{31}{preliminairescomplexite}{Préliminaires}
\INFcinqsixun{32}{exalgorithme}{Exemple d'algorithme}
\INFcinqsixun{33}{algorithme}{Algorithme}
\INFcinqsixun{34}{modeles}{Modèles de calcul}
\INFcinqsixun{35}{calculabilite}{Calculabilité}
\INFcinqsixun{36}{circuitscinqsixun}{Circuits}
\INFcinqsixun{37}{complexitetemps}{Complexité en temps}
\INFcinqsixun{38}{complexiteespace}{Complexité en espace}
\INFcinqsixun{39}{classesparalleles}{Classes Parallèles}
\INFcinqsixun{40}{algoprobabiliste}{Algorithmes probabilistes}
\INFcinqsixun{41}{classesprobabilistes}{Classes probabilistes}
\INFcinqsixun{42}{hierarchiepolynomiale}{Hiérarchie polynomiale}
\INFcinqsixun{43}{preuveinteractive}{Preuves interactives}
\INFcinqsixun{44}{tderandomization}{Derandomization}
\INFcinqsixun{45}{talgoapprox}{Algorithmes d'approximation}
\INFcinqsixun{46}{tquelquepart}{Autres choses quelquepart}
\INFcinqsixun{47}{tvracannexe}{Vrac annexe}


\MPRI{02-B}{diophantine}{Diophantine Equations}

\COURS{10}{circuitsbooleens}{Boolean circuits}
\COURS{02}{BSSmodelearbitrarystructureinshort}{Circuits over a
structure}
\COURS{10}{circuitsoverastructure}{Circuits over a structure}
\COURS{11}{complexiteentemps}{Time complexity}

\COURS{90}{DiscreteODE}{Discrete Ordinary Differential Equations}

%========================================
%                      Poly english
%========================================

\COURS{00}{computationtheorybook}{Computation theory: Logic, Models of
computation, Complexity}

\COURS{00}{polyinfquatrecentdouzeenglish}{Logic, Models of
computation, Complexity}
\COURS{00}{polyenglish}{Logic, Models of
computation, Complexity}
\COURS{05}{premierordreenglish}{Predicate calculus}
\COURS{02}{preliminairesenglish}{Recurrence and induction}
\COURS{38}{complexiteespaceenglish}{Space complexity}
\COURS{39}{problemescompletsenglish}{Complete problems}
\COURS{32}{exalgorithmeenglish}{Examples of algorithms}
\COURS{04}{preuvespropositionnellesenglish}{Proofs}
\COURS{05}{propositionnelenglish}{Propositional calculus}
\COURS{06}{completudeenglish}{Models. Completeness.}
\COURS{08}{algorithmesenglish}{Sequential Algorithms. Church-Turing Thesis. }
\COURS{11}{basecomplexiteenglish}{Basic of complexity analysis of
algorithms} 
\COURS{12}{oracle}{Computing with Oracles. Relativisation}
\COURS{16}{correctionsenglish}{Solutions of some exercises}
\COURS{39}{classesparallelesenglish}{Parallel classes}
\COURS{40}{algoprobabilisteenglish}{Probabilistic algorithms}
\COURS{41}{classesprobabilistesenglish}{Probabilistic classes}
\COURS{42}{hierarchiepolynomialeenglish}{Polynomial hierarchy}
\COURS{43}{preuveinteractiveenglish}{Interactive proofs}
\COURS{44}{tderandomizationenglish}{Derandomization}
\COURS{45}{talgoapproxenglish}{Approximation algorithms}
\COURS{46}{tquelqueparteenglish}{Something somewhere else}
\COURS{47}{tvracannexeenglish}{Vrac appendix}

%========================================
%                      Complexité / Implicite
%========================================

\COURS{31}{coursarnaud}{Théorie des modèles finis et applications}
\COURS{30}{polyLMFI2020}{Descriptive complexity: \\ From the discrete to
the continuum}
\COURS{26}{complexiteimplicite}{Implicit Complexity}
\COURS{27}{rametespace}{RAM machines and complexity classes}
\COURS{46}{boundedarithmetic}{Bounded Arithmetic}
\COURS{47}{potentiel-lfmi}{Potentiel pour LFMI}


%========================================
%                      CSE304
%========================================

\COURS{30}{polyCSE304}{CSE 304. Complexity}

%========================================
%                      Complexité / BSS
%========================================

\MPRI{01}{BSS}{The Blum Shub Smale Model}
% \MPRI{02}{BSSmodele}{The BSS model: definitions}
\MPRI{02}{BSS-circuits}{Circuits and Complexity}
\MPRI{02}{BSSresults}{The BSS model: some other results}


\COURS{02}{BSSmodelearbitrarystructure}{Computing over an arbitrary
structure}
\COURS{01}{BSSmodele-definition}{The BSS model over $\R$}
\COURS{01}{BSSmodelerelationsdiscretecomplexity}{Relation with
discrete complexity}

\COURS{02}{BSSpartialrecursiveness}{A characterization of partial
recursiveness}

\COURS{03}{BSSdescriptivecomplexity}{Descriptive Complexity for BSS machines over the reals}



%====================
%                      Calculabilité
%=====================



%====================
%                      Calculabilité /ITTMS
%=====================

\COURS{20}{ittmeries}{ITTMeries}
\COURS{20}{arithmeticalhyperarithmeticalhierarchy}{Arithmetical and
Hyper-arithmetical hierarchy}
\COURS{21}{ordinalrecursive}{Recursive ordinals}
\COURS{22}{reglelimsup}{The limsup rule}
\COURS{22}{ittmgaps}{Gaps in ordinal computations}
\COURS{23}{ittmpuissance}{Strengh of ITTMs}
\COURS{09}{calculabiliteenglish}{Computability}

%====================
%                      Reverse maths
%=====================

\COURS{28}{reversemath}{Reverse Mathematics}

%====================
%                      Computable Récursive
%=====================

\MPRI{04}{computableanalysis}{Computable Analysis}
\MPRI{05}{computableanalysisbasics}{Computable Analysis}
\MPRI{06}{computableanalysisMPRI}{Computable Analysis}
\MPRI{07}{vraccomputableanalysis}{Computable Analysis: Misc}


\COURS{16}{parameterizedcomplexity}{Parameterized complexity}

\COURS{21}{hierarchiereels}{Hierarchy of Reals}

\COURS{22}{vracanalyserecursive}{Misc. Computable analysis}

%========================================
%                      GPAC / simulations MT
%========================================


\MPRI{02}{Staticundecidability}{Static Undecidability}
\MPRI{02-A}{richardson}{Richardson Theorem}
% ailleurs: \MPRISUB{02-B}{diophantine}{Diophantine Equations}


\MPRI{03}{dynamicundecidability}{Dynamic Undecidability}
\MPRI{04}{basicsODEdynamicalsystems}{Basics about dynamical systems}
\MPRI{05}{pslashpoly}{Non Uniform Polynomial Time}
\MPRI{06}{nonrationalcoefficients}{Non-rationnal coefficients}


\COURS{19}{simulationMTODE}{Simulating Turing Machines by
  Analytic Functions}

%====================
%                      GPAC / bases
%=====================

\COURS{17}{gpac}{The General
  Purpose Analog Computers. Differential Analyzers}


\COURS{18}{pivp}{(GPAC)
  Generable Functions}

%========================================
%                      Calcul formel / classique
%=========================================

\COURS{14}{mathutiles}{Some usefull maths}

\COURS{15}{solvingode}{Solving ODEs}
\COURS{16}{polymathode}{Poly Ordinary Differential Equations}

\COURS{27}{series}{Series}


\COURS{13}{fonctionsanalytiques}{Analytic
  Functions}

\COURS{23}{dfinies}{Holonomous functions. D-finite function}

%========================================
%                      Calcul formel/Equations différentielles
%========================================


\COURS{15}{equationsdifferentielles}{Differential
  equations}


\COURS{15-B}{equationsdifferentielleselucubrations}{Differential
  equations - Elucubrations}

%========================================
%                      Calcul formel / surréels+transséries
%=========================================

\COURS{25}{surreels}{Surreals
  and co}
\COURS{25}{surrelsfrancais}{Surreals
  and co}
\COURS{25}{surreelselucubrations}{Surreals
  and co}


%========================================
%                      Misc.
%=========================================

\COURS{99}{therest}{Misc. The rest.}

